\chapter{Related Works}
\label{ch:related}

This chapter introduces the algorithms used, the reasoning behind the evaluation metric
chosen, and the statistical procedures on which the comparison rests.

\section{Classifiers for Tabular Fraud Detection}

Early automated fraud detection relied on linear models and simple decision trees.
\citet{dalpozzolo2015calibrating} studied credit card fraud detection in detail and showed
that the way data is sampled and evaluated has a large effect on reported results, a
finding this project also observes.

Logistic Regression predicts fraud by fitting a straight decision boundary through
the feature space. It is fast and easy to interpret, but it cannot represent interactions
between features unless those interactions are built in by hand.

Random Forest \citep{breiman2001random} averages many decision trees, each fitted
to a bootstrap sample of the rows and a random subset of the features at each split. The
averaging reduces variance, and the randomisation keeps the trees from all making the same
mistakes.

Gradient boosting fits trees one after another, with each new tree correcting the
mistakes of the trees built so far. XGBoost \citep{chen2016xgboost} adds an explicit
penalty on tree complexity to this process, which helps control overfitting.

LightGBM \citep{ke2017lightgbm} is the model this project ultimately deploys. It
differs from XGBoost in three ways that matter on a dataset of this size. It grows trees
leaf-wise, meaning it always splits the leaf with the largest loss reduction rather than
expanding every node at one depth, which reaches a lower loss for the same number of
leaves. It uses \gls{goss}, keeping all the instances with large gradients while randomly
subsampling the well-fitted ones, which suits imbalanced data. Finally, \gls{efb} merges
mutually exclusive sparse features to reduce dimensionality, and histogram-based split
finding makes training fast even on hundreds of thousands of rows.

\section{Strategies for Class Imbalance}

Four approaches are compared in this report.

Class weighting leaves the data untouched and instead makes each fraud count for
more when the model is trained. The weights come from the ``balanced'' heuristic of
scikit-learn \citep{pedregosa2011scikit}, which sets each class weight inversely
proportional to that class's frequency, so on this training set a fraud counts for roughly
599 times as much as a legitimate transaction.

Random undersampling discards majority-class rows until the classes are balanced.
It is very simple, but on data this imbalanced it is destructive: balancing 284 frauds
against 169{,}951 legitimate rows means throwing away 169{,}667 of them.

Random oversampling duplicates minority rows instead. No information is discarded,
but exact duplicates give the model the same evidence several times over, which can
encourage it to overfit the particular frauds it has seen.

SMOTE \citep{chawla2002smote} creates new synthetic minority examples by
interpolating between existing fraud cases and their nearest fraud neighbours, rather than
copying them outright.

All three resampling methods share a serious pitfall. If they are applied to the whole
dataset before splitting, synthetic or duplicated rows derived from test-set frauds leak
into training, and the reported performance becomes fiction. For this reason every
resampler in this project is a step inside the pipeline, so it is fitted only on the
training fold.

\section{Evaluation under Severe Imbalance}

Let \gls{tp}, \gls{fp}, \gls{fn}, and \gls{tn} denote the four confusion-matrix counts. The
two metrics used throughout this report are then
\begin{equation}
  \label{eq:precisionrecall}
  \mathrm{Precision} = \frac{TP}{TP+FP}, \qquad
  \mathrm{Recall} = \frac{TP}{TP+FN}.
\end{equation}
Equation~\ref{eq:precisionrecall} gives the standard definitions of precision and recall
from the information retrieval literature, in the form used by
\citet{davis2006relationship} in their analysis of ROC and precision-recall curves. Neither
is original to this project.

Precision is the share of flagged transactions that really are fraud, and recall is the
share of real frauds that the model flags. The $F_1$ score reported in later tables is
their harmonic mean, which is high only when both are high.

Accuracy, $(TP+TN)/n$, is the standard textbook definition and is unusable here, because
the trivial all-legitimate classifier
scores 0.9983. \gls{rocauc} is also optimistic, because the enormous pool of true negatives
barely moves it even when thousands of false alarms occur. The precision-recall curve is
more honest about the errors that matter, because precision is sensitive to every false
alarm \citep{davis2006relationship,saito2015precision}. This project therefore ranks models
by PR-AUC, the area under this curve, computed as average precision, which is the
recall-weighted mean of precision across every possible decision threshold. A random
classifier scores the prevalence of the positive class, which here is 0.00167, rather than
the 0.5 that ROC-AUC would give.

Two further metrics are also used. The \gls{mcc} \citep{matthews1975comparison} summarises
all four confusion-matrix cells into a single score that stays informative under imbalance.
The Brier score \citep{brier1950verification} measures how well a predicted probability
matches the true outcome, and pairs with the calibration analysis.

\section{Statistical Comparison of Classifiers}
\label{sec:statmethods}

Comparing two classifiers on the same test set is a paired problem: the two models see
identical transactions, so their errors are correlated and a plain unpaired test would be
invalid \citep{dietterich1998approximate}. McNemar's test \citep{mcnemar1947note} is the
appropriate paired test. It looks only at the cases where the two models disagree, and asks
whether one model wins those disagreements more often than chance would allow. Because all
28 model pairs are tested here, the resulting $p$-values are adjusted
using the Holm step-down procedure \citep{holm1979simple}, which corrects for running many
comparisons at once.

For interval estimates around PR-AUC, a percentile bootstrap \citep{efron1987better} is
used: the test set is resampled with replacement 1{,}000 times and the metric recomputed
each time, giving a range of plausible values rather than a single number. For precision
and recall, the Wilson score interval \citep{wilson1927probable} is used instead, since it
stays within $[0,1]$ and remains accurate even for the small counts involved here.
